\documentclass[11pt]{article}
\usepackage[T1]{fontenc}
\usepackage[margin=1in]{geometry}
\usepackage{amsmath,amssymb,amsthm,hyperref}
\newtheorem{theorem}{Theorem}
\newtheorem{lemma}[theorem]{Lemma}
\newtheorem{corollary}[theorem]{Corollary}
\DeclareMathOperator{\Aut}{Aut}
\DeclareMathOperator{\Cay}{Cay}
\DeclareMathOperator{\Stab}{Stab}
\title{Complete non-normal orbit quotients of every order at least three}
\author{Main-agent research note for the math-open-problems repository}
\date{17 September 2026}
\begin{document}
\maketitle
\begin{abstract}
For every integer $m\geq3$, we construct by an existence theorem an infinite,
connected, locally finite, simple transitive graph $X$ with a freely acting,
non-normal subgroup $B\leq\Aut(X)$ such that $X_B\cong K_m$, while no normal
subgroup of any transitive subgroup of $\Aut(X)$ has an isomorphic orbit
quotient. The key input is the generating-set-preserving graphical regular
representation construction of Leemann and de la Salle. This extends the
repository's earlier three-vertex candidate. It addresses the literal
formulation of Grimmett--Li Question 4.5 using finite quotient graphs, not an
additional infinite-quotient requirement. Independent audits and publication
novelty review remain pending.
\end{abstract}

\section{Source, attribution, and scope}
Grimmett and Li \cite[Section 4.2, Question 4.5]{GL} ask whether a non-normal
subgroup of a transitive automorphism group can have a transitive simple
orbit quotient that is not obtainable from a normal subgroup of any
transitive automorphism subgroup of the same graph. They additionally ask
about free actions. We use exactly their simple undirected orbit convention:
distinct orbits are adjacent when some representatives are adjacent.
Loops and parallel edges are suppressed.

The ambient graph below is infinite. Its quotient is finite. The displayed
Question 4.5 does not impose infinitude of the quotient. Thus the result gives
a YES answer with that literal scope; it does not settle the variant where
the quotient must also be infinite, nor their separate Question 4.6 about
connective constants.

The first repository candidate is P20-D0072-NONNORMAL-QUOTIENT-WITNESS-01,
using $S_3*A_5$ and $K_3$, in the P20 report dated
2026-09-17T16-18-23Z, creation commit
\texttt{92b25a56025b3e0905274c554a4104d4acbc0685}.
This note credits that candidate and supplies a uniform extension; it does
not claim to have originated its GRR mechanism. Classical finite-group facts
used below include generation of alternating groups by three-cycles and
simplicity of $A_m$ for $m\geq5$; see \cite[Chapter 4]{Milne}.

\section{The published input}
\begin{lemma}[Relative GRR input]
Let $\Lambda$ be a finitely generated group that is not virtually abelian,
and let $S_0\subseteq\Lambda\setminus\{1\}$ be a finite symmetric generating
set. There is a finite symmetric set $S\supseteq S_0$, not containing the
identity, such that $\Aut(\Cay(\Lambda,S))=\Lambda_L$, the left-regular action.
\end{lemma}
\begin{proof}
This is the construction in the proof of Theorem 1.1 of \cite{LdS}, not
merely a claim about its abstract. That proof begins with an arbitrary
$S_0$, sets $S_1=(S_0\cup S_0^2\cup S_0^3)\setminus\{1\}$, and applies
Proposition 3.2 to obtain $S\supseteq S_1$. Its automorphisms preserve the
$S_1$-colours; Theorem 3.1 then makes each one a left translation.
In particular, all prescribed edges from $S_0$ remain present.
\end{proof}

\section{A finite-template criterion}
\begin{lemma}\label{lem:template}
Suppose a finite group $F$ has a faithful transitive action on a set of
$m\geq3$ points, has a nontrivial point stabilizer $H$, and every homomorphism
from $F$ to any group of order $m$ is trivial. Then there is an infinite,
connected, locally finite, simple transitive graph $X$ and a non-normal,
freely acting subgroup $B\leq\Aut(X)$ with $X_B\cong K_m$, but
$X_N\not\cong K_m$ whenever $N\lhd\Gamma'$ and
$\Gamma'\leq\Aut(X)$ acts transitively.
\end{lemma}
\begin{proof}
Take $\Lambda=F^{(1)}*F^{(2)}$, the free product of two copies of $F$.
Let $\pi:\Lambda\to F$ be the identity on the first factor and trivial on
the second, and put $B=\pi^{-1}(H)$. Since $\pi$ is onto,
$[\Lambda:B]=[F:H]=m$.

We verify the hypotheses of the relative GRR input. The group $\Lambda$ is
finitely generated. Choose distinct nonidentity elements $a,b$ of the first
factor and a nonidentity element $c$ of the second. There are $2^k$ distinct
reduced words $x_1c\cdots x_kc$ with $x_i\in\{a,b\}$, all of length $2k$
in the generating set consisting of the nonidentity elements of the two
factors. The free-product normal form proves their distinctness, so
$\Lambda$ is infinite and has exponential growth. A finitely generated
virtually abelian group has polynomial growth: pass to a finite-index
normal abelian subgroup, choose finitely many coset representatives, and
rewrite a word of length $r$ as one representative followed by at most $r$
factors in a fixed finite set of the abelian subgroup. Their exponent
vectors have at most polynomially many possibilities. Thus $\Lambda$ is
not virtually abelian.

Set
\[
 S_0=(F^{(1)}\setminus\{1\})\cup(F^{(2)}\setminus\{1\}).
\]
Choose $S\supseteq S_0$ by the relative GRR input, and let
$X=\Cay(\Lambda,S)$, with an edge from $g$ to $gs$ for $s\in S$.
The graph is infinite and simple because $1\notin S$, locally finite
because $S$ is finite, connected because $S$ generates, and transitive
under left translations. Its full automorphism group is $\Lambda_L$.
We identify $B$ with the corresponding subgroup of left translations.

The point stabilizer $H$ is non-normal. Indeed, if it were normal, an
element fixing the chosen point would, by conjugation and transitivity,
fix every point. Faithfulness would imply $H=\{1\}$, contrary to the
hypothesis. Surjectivity of $\pi$ now implies that $B$ is non-normal in
$\Lambda$. Its vertex action is free: $bg=g$ implies $b=1$.

The $B$-orbits are $Bg$, which are \emph{right cosets} despite arising from
left multiplication. The map $Bg\mapsto H\pi(g)$ is a bijection onto the
$m$ right cosets of $H$ in $F$. To see the injective part explicitly,
$H\pi(g)=H\pi(g')$ implies $\pi(g'g^{-1})\in H$, and hence
$g'g^{-1}\in B$, so $Bg=Bg'$. Surjectivity follows from that of $\pi$.
For distinct cosets choose representatives $r,t$ in the first copy of $F$.
Then $r^{-1}t\neq1$ belongs to $S_0$, so $r$ and $t$ are adjacent in $X$.
Every two distinct quotient vertices are therefore adjacent. The simple
orbit convention gives $X_B\cong K_m$.

It remains to handle \emph{every} transitive automorphism subgroup, not
only a nominated one. If $\Gamma'\leq\Lambda_L$ is transitive, its orbit
of the identity contains every element of $\Lambda$. Regularity means
that the unique left translation sending $1$ to $g$ must belong to
$\Gamma'$ for each $g$. Thus $\Gamma'=\Lambda_L$.

Suppose $N\lhd\Gamma'$ and $X_N\cong K_m$. Under this regular
identification, $N$ corresponds to a normal subgroup $N_0\lhd\Lambda$,
and the number of quotient vertices is $[\Lambda:N_0]=m$.
Consequently $Q=\Lambda/N_0$ is a group of order $m$.
The quotient homomorphism restricted to either free factor is trivial by
the hypothesis on $F$. Since these factors generate $\Lambda$, the whole
quotient homomorphism is trivial. It cannot surject onto a group of order
$m\geq3$. This contradiction proves the asserted universal exclusion.
\end{proof}

\section{All quotient orders $m\geq3$}
\begin{theorem}\label{thm:all}
For every integer $m\geq3$, the conclusion of Lemma~\ref{lem:template}
holds. In particular, the literal Grimmett--Li Question 4.5 has a YES
answer even when the non-normal subgroup must act freely.
\end{theorem}
\begin{proof}
Use the natural faithful transitive action on $m$ letters and take
\[
 F=\begin{cases}S_3,&m=3,\\ A_m,&m\geq4.\end{cases}
\]
The point stabilizer has order $2$ for $m=3$ and $(m-1)!/2$ for $m\geq4$,
so is nontrivial in every case. Transitivity is immediate for $S_3$;
for $A_m$, a three-cycle sends any chosen letter to any other, using a
third letter.

Let $Q$ be any group of order $m$. For $m=3$, every transposition in
$S_3$ must map to the identity, since $Q$ has no element of order $2$.
Transpositions generate $S_3$, so every homomorphism is trivial.
For $m=4$, every three-cycle in $A_4$ must map to the identity, since
$Q$ has no element of order $3$. Three-cycles generate $A_4$, so again
all homomorphisms are trivial.
For $m\geq5$, simplicity of $A_m$ makes any nontrivial homomorphism
injective. But $|A_m|=m!/2>m=|Q|$, which rules out injectivity.
Thus the last hypothesis of Lemma~\ref{lem:template} holds in all cases,
and that lemma gives the result.
\end{proof}

\begin{corollary}[An explicit, non-optimal degree bound]
Write $f=|F|$, $a=f-1$, and $t=2(a+a^2+a^3)$. The graph supplied above
may be chosen of degree at most $2t(t+14)$.
\end{corollary}
\begin{proof}
By the free-product normal form there are exactly $2a^j$ reduced words
of length $j\geq1$ in the full factor generating set $S_0$.
Consequently $|S_1|=2(a+a^2+a^3)=t$.
Proposition 3.2 of \cite{LdS} gives $|S|\leq2|S_1|(|S_1|+14)$.
Since the graph is a simple Cayley graph with $1\notin S$, its degree
is $|S|$, proving the bound.
\end{proof}
For $m=3$, this bound is $200880$; it is not an optimality claim and it
does not amount to listing the final generators of a GRR.

\section{Verification, failed route, and limitations}
The accompanying standard-library Python checker uses exact permutation
composition. It checks the $S_3,A_4,A_5$ templates: group orders, point
stabilizers, non-normality, coset partitions, complete quotient adjacency,
and the entire normal-subgroup lattice of each finite group. It verifies
that a homomorphic image whose order divides $m$ must be trivial. The
normal-subgroup search is exhaustive because every normal subgroup is
a union of conjugacy classes and is reached by successively adjoining
those classes. The archived run ends \texttt{ALL\_EXACT\_CHECKS\_PASS}.
This is a self-check, not an independent audit or a machine verification
of the infinite GRR input or of all alternating groups.

A tempting replacement $F=S_m$ for all $m$ does not support this proof:
when $m$ is even, the sign maps allow $S_m*S_m$ to map onto a group of
order $m$ generated by two involutions (dihedral, with the order-four
case given by the Klein group). Thus abelianization alone does not
exclude normal index-$m$ subgroups. The alternating-group choice avoids
this obstruction; it does not follow from a computation on a few values
of $m$.

The source statement and the relative-GRR proof were checked in their
primary PDFs on 17 September 2026. Bounded searches for the exact question,
non-normal quotients with GRRs, and complete orbit quotients did not locate
an explicit prior resolution. That is not a proof of novelty or continuing
worldwide openness. Related Schreier-graph literature requires further
review. The GRR theorem and the finite-group facts are published inputs,
not new results here. No claim is made to solve Sidorenko, graph
reconstruction, or another headline conjecture. This note is not campaign
closure until the required two distinct OTHER-track audits are recorded;
this run creates no such audit votes.

\begin{thebibliography}{9}
\bibitem{GL} G. R. Grimmett and Z. Li,
\emph{Self-avoiding walks and connective constants},
arXiv:1704.05884v2, 2019, Section 4.2, Question 4.5, pp. 16--17.
\url{https://arxiv.org/abs/1704.05884}.
\bibitem{LdS} P.-H. Leemann and M. de la Salle,
\emph{Cayley graphs with few automorphisms: the case of infinite groups},
Annales Henri Lebesgue \textbf{5} (2022), 73--92.
Theorem 1.1, Theorem 3.1, Proposition 3.2, and Section 5.
\url{https://doi.org/10.5802/ahl.118}.
\bibitem{Milne} J. S. Milne, \emph{Group Theory}, version 4.01 (2025),
Chapter 4, permutation groups.
\url{https://www.jmilne.org/math/CourseNotes/GT.pdf}.
\end{thebibliography}
\end{document}
